% ---------------------------------------------------------------------
\section{标度无关性与 DIS 方案}\label{sec:disall-scale}

$\mu$ 是把共线发散分给 PDF 时引入的人为标度。DGLAP 方程（\cref{thm:pdf-dglap}）是算符重整化的结论，逻辑上先于因子化定理；标度无关性是它与系数函数的 $\mu$ 依赖相互抵消的推论。

\begin{proposition}[标度无关性 \tagR]\label{prop:disall-scale}
（$R_{\mathrm{PT}}$：在重整化微扰论中对红外正规化的部分子态严格。）以下陈述针对截断到 $\order{\as}$ 的\cref{eq:disall-norm}。（本命题只涉及部分子态的系数函数。）
设 PDF 满足 DGLAP 方程，$a\equiv\as(\mu_R)/2\pi$ 满足 $\dd a/\dd\ln\mu_R^2=-\frac{\beta_0}2a^2+\order{a^3}$，系数函数取\cref{prop:disall-nlo}的 $C_2^i(z,Q^2/\mu_F^2)$。则
\begin{equation}\label{eq:disall-scale}
\frac{\partial}{\partial\ln\mu_F^2}\frac{F_2}x=\order{a^2},\qquad
\frac{\partial}{\partial\ln\mu_R^2}\frac{F_2}x=-\frac{\beta_0}2a^2\sum_qe_q^2\big[C_2^q\otimes(q+\bar q)+2C_2^g\otimes g\big]+\order{a^3},   % VF-DIS-10 (written for F_2/x)
\end{equation}
式中 $C_2^i$ 即 $C_2^i(z,Q^2/\mu_F^2)$；且 $F_L$ 在 $\order{a}$ 与 $\mu_F$、$\mu_R$ 都无关。
\uses{prop:disall-nlo,thm:pdf-dglap,prop:prelim-beta}
\end{proposition}

\begin{proof}
在 Mellin 空间中卷积成为乘积。对一味夸克，记 $\Sigma_N$、$g_N$ 为 $q+\bar q$ 与 $g$ 的矩，$F_{2,N}$ 为 $F_2/x$ 的矩，$L_F=\ln(Q^2/\mu_F^2)$，则
$F_{2,N}=e_q^2\{\Sigma_N(\mu_F)+a(\mu_R)[(C^q_{2,N}+L_FP_{qq,N})\Sigma_N+2(C^g_{2,N}+L_FP_{qg,N})g_N]\}$。由\cref{thm:pdf-dglap}，
$\partial\Sigma_N/\partial\ln\mu_F^2=a(\mu_F)[P_{qq,N}\Sigma_N+2P_{qg,N}g_N]+\order{a^2}$（$q$ 与 $\bar q$ 各从胶子得到一份 $P_{qg}\otimes g$），而 $\partial L_F/\partial\ln\mu_F^2=-1$，故
$\partial F_{2,N}/\partial\ln\mu_F^2=e_q^2[a(\mu_F)-a(\mu_R)](P_{qq,N}\Sigma_N+2P_{qg,N}g_N)+\order{a^2}=\order{a^2}$。对 $\mu_R$，只有 $a(\mu_R)$ 依赖于它，
由\cref{prop:prelim-beta}，$\dd a/\dd\ln\mu^2=\frac1{2\pi}\dd\as/\dd\ln\mu^2=-\frac{\beta_0}{8\pi^2}\as^2=-\frac{\beta_0}2a^2$，即得第二式。
$F_L$ 在 $\order{a}$ 只含与 $\mu$ 无关的 $C_L$ 乘以 PDF，而 PDF 的 $\mu$ 依赖是 $\order{a}$，故 $F_L$ 的变化为 $\order{a^2}$。
（以上用符号计算在 Mellin 空间中核对过。）
\end{proof}

\noindent
定理只保证剩余的标度依赖是高一阶的；把 $\mu$ 在 $Q/2$ 与 $2Q$ 之间变化的差别当作“理论误差”是没有定理支持的约定。
另一种常用方案把 $F_2$ 的全部高阶修正吸收进夸克分布；方案变换的一般形式是 $f'=Z\otimes f$、$C'=C\otimes Z^{-1}$（\cref{rem:pdf-schemes}）。

\begin{proposition}[DIS 方案 \tagR]\label{prop:disall-dis-scheme}
（$R_{\mathrm{PT}}$：在重整化微扰论中对红外正规化的部分子态严格。）以下陈述在 $\order{\as}$。（本命题只涉及部分子态的系数函数与方案变换。）
DIS 方案要求在 $\mu=Q$ 处 $F_2/x=\sum_qe_q^2(q+\bar q)^{\mathrm{DIS}}$ 对各阶成立。在 $\order{\as}$（$a=\as/2\pi$，全部量取在同一 $\mu$，$C_2^i\equiv C_2^i(z)$ 为 $\mu=Q$ 处的 $\MSbar$ 系数），
\begin{equation}\label{eq:disall-dis-scheme}
\begin{gathered}
q_i^{\mathrm{DIS}}=q_i+a\big[C_2^q\otimes q_i+C_2^g\otimes g\big],\qquad \bar q_i^{\mathrm{DIS}}=\bar q_i+a\big[C_2^q\otimes\bar q_i+C_2^g\otimes g\big],\\
g^{\mathrm{DIS}}=g-a\big[C_2^q\otimes\Sigma+2n_fC_2^g\otimes g\big],   % VF-DIS-11
\end{gathered}
\end{equation}
$\Sigma=\sum_i(q_i+\bar q_i)$；胶子的变换不由 $F_2$ 确定，上式是对全部矩保持动量求和规则的选择。于是 $C_2^{q,\mathrm{DIS}}=\ln(Q^2/\mu^2)P_{qq}$，$C_2^{g,\mathrm{DIS}}=\ln(Q^2/\mu^2)P_{qg}$；
领头阶分裂函数不变，次领头阶在 $(\Sigma,g)$ 基中按
\begin{equation}\label{eq:disall-P1}
\begin{gathered}
P'^{(1)}=P^{(1)}+Z\otimes P^{(0)}-P^{(0)}\otimes Z-\frac{\beta_0}2Z,\qquad Z=\begin{pmatrix}C_2^q&2n_fC_2^g\\-C_2^q&-2n_fC_2^g\end{pmatrix},\\
P^{(1),\mathrm{DIS}}_{\mathrm{NS}}=P^{(1),\MSbar}_{\mathrm{NS}}-\frac{\beta_0}2C_2^q   % VF-DIS-11
\end{gathered}
\end{equation}
变换（DGLAP 核写作 $aP^{(0)}+a^2P^{(1)}$）；$C_L$ 在 $\order{\as}$ 与方案无关。
\uses{prop:disall-nlo,rem:pdf-schemes,thm:pdf-dglap,prop:prelim-beta}
\end{proposition}

\begin{proof}
(1) 由\cref{eq:disall-dis-scheme}，$q_i=q_i^{\mathrm{DIS}}-a[C_2^q\otimes q_i+C_2^g\otimes g]+\order{a^2}$。代入\cref{eq:disall-norm}（取一般 $\mu$，$C_2^i(z,Q^2/\mu^2)=C_2^i(z)+\ln(Q^2/\mu^2)P_{qi}$），
$\order{a}$ 中的 $C_2^q\otimes(q+\bar q)+2C_2^g\otimes g$ 恰好抵消，剩下
$F_2/x=\sum_qe_q^2\{(q+\bar q)^{\mathrm{DIS}}+a\ln(Q^2/\mu^2)[P_{qq}\otimes(q+\bar q)+2P_{qg}\otimes g]\}+\order{a^2}$，即 $C_2^{i,\mathrm{DIS}}=\ln(Q^2/\mu^2)P_{qi}$，在 $\mu=Q$ 处为零。
(2) $\Sigma^{\mathrm{DIS}}+g^{\mathrm{DIS}}=\Sigma+g+a[C_2^q\otimes\Sigma+2n_fC_2^g\otimes g]-a[C_2^q\otimes\Sigma+2n_fC_2^g\otimes g]=\Sigma+g$，对每个矩都成立，特别地 $N=2$ 的动量求和规则不变。
(3) $F_L$ 是 $\order{a}$ 的量，PDF 的 $\order{a}$ 变化只改变 $\order{a^2}$ 项，所以 $C_L$ 不变；领头阶分裂函数也不变。
(4) 设 $f'=(1+aZ)\otimes f$，$\dd f/\dd\ln\mu^2=(aP^{(0)}+a^2P^{(1)})\otimes f$，$\dd a/\dd\ln\mu^2=-\frac{\beta_0}2a^2$（\cref{thm:pdf-dglap,prop:prelim-beta}）。则
$\dd f'/\dd\ln\mu^2=(1+aZ)\otimes(aP^{(0)}+a^2P^{(1)})\otimes f-\frac{\beta_0}2a^2Z\otimes f$；用 $f=(1-aZ)\otimes f'+\order{a^2}$ 消去 $f$，得
$\dd f'/\dd\ln\mu^2=aP^{(0)}\otimes f'+a^2[P^{(1)}+Z\otimes P^{(0)}-P^{(0)}\otimes Z-\frac{\beta_0}2Z]\otimes f'+\order{a^3}$。在 $(\Sigma,g)$ 基中 $Z$ 的两行分别来自 $\Sigma^{\mathrm{DIS}}$ 与 $g^{\mathrm{DIS}}$ 的定义。
非单态部分 $Z_{\mathrm{NS}}=C_2^q$ 在 Mellin 空间中与 $P^{(0)}_{\mathrm{NS}}$ 对易，于是 $P^{(1),\mathrm{DIS}}_{\mathrm{NS}}=P^{(1),\MSbar}_{\mathrm{NS}}-\frac{\beta_0}2C_2^q$。
（以上步骤均以符号计算核对。）
\end{proof}

\noindent
夸克分布的 DIS 方案可以追溯到\cite{AltarelliEllisMartinelli1979}；用动量求和规则固定胶子的约定参见\cite{DiemozEtAl1988}。其他过程（例如 Drell--Yan）的硬系数必须按同一个 $Z$ 变换，
在那里 DIS 方案的系数并不为零；方案的选择同时改变 PDF 与系数函数，可观测量只在所截断的阶以上改变。
